\documentclass{article}
\usepackage{amsmath,amssymb,amsthm}
\newtheorem{claim}{Claim}
\title{Pure-mathematics lemmas behind the PhaseTripleProduct (PTP) resource and control theorems}
\author{Draft note for independent audit and formalization}
\date{}
\begin{document}
\maketitle

\section*{Scope}
This note states six self-contained mathematical claims used by a verified resource analysis of a
recursive Toom--Cook phase-triple-product circuit. Each claim is elementary and independent of any
circuit model. For each we give a proof sketch. Nothing in the claims refers to quantum gates: a claim is
a statement about integers, rationals, reals and finite lists. (The gate-level consequences are
proved separately in Lean and are not claimed here.)

\section{Two's-complement masking}
For $n\ge 1$ let $\mathrm{tc}_n:\{0,\dots,2^n-1\}\to\mathbb Z$ be
$\mathrm{tc}_n(m)=m$ if $m<2^{n-1}$ and $m-2^n$ otherwise.

\begin{claim}\label{c:mask}
For $c\in\{0,1\}$ and $0\le m<2^n$, $\mathrm{tc}_n(c\cdot m)=c\cdot \mathrm{tc}_n(m)$.
Moreover, if $m'$ is obtained from $m$ by replacing every binary digit $d_i$ by $c\wedge d_i$, then $m'=c\cdot m$, so $\mathrm{tc}_n(m')=c\cdot\mathrm{tc}_n(m)$.
\end{claim}
\begin{proof} If $c=1$ every digit is unchanged, so $m'=m$ and the identity is trivial. If $c=0$ every digit becomes $0$, so $m'=0$, and both sides are $\mathrm{tc}_n(0)=0$ (this uses $n\ge1$, so that the threshold $2^{n-1}$ is defined).\end{proof}

\section{Linear growth of the shift/add width scan}
Fix an integer $k\ge 2$. A \emph{width state} is a bound on magnitudes ($|\text{value}|<2^{w_i}$, so negation keeps $w_i$; it is not a two's-complement width), $w:\{0,\dots,k-1\}\to\mathbb N$. A \emph{program} is a finite list of
operations, with transitions
\begin{align*}
\mathrm{shiftL}(i,s)&: w_i \mapsto w_i+s, &
\mathrm{shiftR}(i,s)&: w_i\mapsto w_i\dot- s,\\
\mathrm{negate}(i)&: w\mapsto w, &
\mathrm{addScaled}(d,e,sh)&: w_d\mapsto 1+\max(w_d,\,w_e+sh),\\
\mathrm{phase}(i)&: w\mapsto w.
\end{align*}
Write $\mathrm{mx}(w)=\max_i w_i$ and let $\mathrm{scan}(w,\text{ops})$ be the maximum of $\mathrm{mx}$ over all
states visited, the initial state included. Define $D(\text{ops})=\sum_{\mathrm{shiftL}(i,s)}s+\sum_{\mathrm{addScaled}(d,e,sh)}(1+sh)$.

\begin{claim}\label{c:scan}
$\mathrm{scan}(w,\text{ops})\le \mathrm{mx}(w)+D(\text{ops})$.
\end{claim}
\begin{proof} Induction on the program. Each transition raises $\mathrm{mx}$ by at most the term it contributes to $D$:
$\mathrm{shiftL}$ by $s$; $\mathrm{addScaled}$ gives $1+\max(w_d,w_e+sh)\le 1+sh+\mathrm{mx}(w)$; the others do not raise it.\end{proof}

\begin{claim}\label{c:child}
Let $n\ge 1$, $W=\lfloor n/k\rfloor$, and let the initial state be $w_j=W$ for $j<k-1$ and $w_{k-1}=n-(k-1)W$.
Then $\mathrm{mx}(w)\le W+k-1$, hence the child width $1+\mathrm{scan}(w,\text{ops})$ is at most $\lfloor n/k\rfloor+k+D(\text{ops})\le\lceil n/k\rceil+c$ with $c=k+D(\text{ops})$.
\end{claim}
\begin{proof} Write $n=kW+r$ with $0\le r<k$; then $n-(k-1)W=W+r\le W+k-1$. Combine with Claim~\ref{c:scan}.\end{proof}

\section{Arity-three Toom--Cook identity}
Let $k\ge2$ be an integer and $q=3(k-1)+1=3k-2$. A \emph{point} is a representative $p=(a,b)\in\mathbb Q^2\setminus\{0\}$ of a point of the projective
line; two points are \emph{distinct} if they are not proportional. (The affine point $a/b$ with $b=1$ and the point at infinity
$(1,0)$ are both allowed.) For a limb vector $x=(x_0,\dots,x_{k-1})$ define the binary form
$X(a,b)=\sum_{j<k}x_j\,a^j b^{\,k-1-j}$, homogeneous of degree $k-1$, and put $x(p)=X(a,b)$ for $p=(a,b)$.
For $B\in\mathbb Q$ the value $X(B,1)=\sum_j x_jB^j$ is the number represented by the limbs at radix $B$.

\begin{claim}\label{c:node}
Let $p_0,\dots,p_{q-1}$ be pairwise distinct (non-proportional, rational) points, with fixed representatives; the $c_t$ depend on the chosen representatives. There are rationals $c_0,\dots,c_{q-1}$ depending only on the points and
on $B$, not on the limbs, such that for all limb vectors $x^{(0)},x^{(1)},x^{(2)}$
\[
X^{(0)}(B,1)\,X^{(1)}(B,1)\,X^{(2)}(B,1)\;=\;\sum_{t=0}^{q-1} c_t\prod_{s=0}^{2}x^{(s)}(p_t).
\]
\end{claim}
\begin{proof} $R=X^{(0)}X^{(1)}X^{(2)}$ is a binary form of degree $3(k-1)=q-1$. The evaluation map
$R\mapsto (R(p_t))_{t<q}$ from binary forms of degree $q-1$ to $\mathbb Q^q$ is injective, because a nonzero binary form of
degree $q-1$ has at most $q-1$ zeros on the projective line; it is therefore bijective by dimension count. Let $c_t$ be the
coordinates of the linear functional $R\mapsto R(B,1)$ in the basis dual to evaluation. Then $R(B,1)=\sum_tc_tR(p_t)$, and
$R(p_t)=\prod_sX^{(s)}(p_t)$ because evaluation is multiplicative.\end{proof}

\section{A master-theorem bound and the exponent}
Let $k\ge2$, $q=3k-2$ and $\alpha=\log q/\log k$, so $k^\alpha=q$.

\begin{claim}\label{c:alpha}
$1<\alpha\le2$, $\alpha_2=2$, and for every $\varepsilon>0$ there is $k$ with $\alpha_k\le 1+\varepsilon$.
\end{claim}
\begin{proof} $q>k$ gives $\alpha>1$. $q\le k^2$ because $k^2-3k+2=(k-1)(k-2)\ge0$, with equality at $k=2$, so $\alpha\le2$ and
$\alpha_2=2$. Also $q<3k$ gives $\alpha<1+\log3/\log k\to1$.\end{proof}

\begin{claim}\label{c:master}
Let $c,A,B\in\mathbb N$ and let $f:\mathbb N\to\mathbb N$ be nondecreasing. Put $m(n)=\lceil n/k\rceil+c$. Suppose there is $n_1$ such
that for all $n\ge n_1$ we have $m(n)<n$ and $f(n)\le An+B+q\,f(m(n))$. Then there is a real $C>0$ with
$f(n)\le C\,n^{\alpha}$ for all $n\ge1$.
\end{claim}
\begin{proof} Sketch. Put $\beta=q/k>1$ and $\Phi(n)=Cn^\alpha-En$ with $E=E(C)\ge 0$ to be chosen. For $n\ge n_0$ we have
$n/k\le m(n)\le n/k+c+1$, hence $m^\alpha\le (n/k)^\alpha+a\,n^{\alpha-1}$ for a constant $a=a(k,c,\alpha)$. Using
$q(n/k)^\alpha=n^\alpha$, $qm\ge\beta n$ and $n^{\alpha-1}\le n$ (as $\alpha\le2$),
\[
An+B+q\,\Phi(m)\le \Phi(n)+n\bigl(Cqa+A+B-E(\beta-1)\bigr)\le\Phi(n)\qquad\text{if }E=\tfrac{Cqa+A+B}{\beta-1}.
\]
Note $E/C\le (qa+A+B)/(\beta-1)$ for $C\ge1$, but $E/C$ may exceed $1$, so $\Phi$ is negative for small $n$ and one cannot
compare $f$ with $\Phi$ from $n=1$. Instead fix $N\ge n_0$ with $m(n)<N$ only for $n$ in a finite range, and run the strong induction
for $n\ge N$: if $m(n)\ge N$ use the induction hypothesis $f(m(n))\le\Phi(m(n))$; if $m(n)<N$ use $f(m(n))\le f(N)$ (monotonicity),
which is a constant, and require $\Phi(n)\ge An+B+q f(N)$ on that finite range, true once $C$ is large. For $n<N$ absorb $f(n)\le f(N)$
into $C$.\end{proof}

\section*{Intended use}
Claim~\ref{c:child} gives the one-level width contraction $n\mapsto \lceil n/k\rceil+c$; Claim~\ref{c:master} then turns a
per-level cost $An+B$ plus $3k-2$ children into $O(n^{\log_k(3k-2)})$ (Claim~\ref{c:alpha});
Claim~\ref{c:node} is the algebraic correctness of one level; Claim~\ref{c:mask} makes a controlled version of the whole
circuit cost only $O(n)$ extra gates.
\end{document}
